%*******************************************************************************
%*********************************** First Chapter *****************************
%*******************************************************************************
\graphicspath{{Figures}}
\epigraphhead[550]{""AI is the science of making machines do things that would require intelligence if done by men."
\par\hfill\textsc{― Marvin Minsky}}
\part{Introduction}  %Title of the First Chapter

% Quote by Marvin Minsky in the book - Semantic Information Processing, MIT Press 1968


% \cleardoublepage
\chapter{Introduction}
\label{ch:intro}
\textit{}
% \minitoc



\section{Context and Motivation}
\label{sec:motivation}
\textit{\textbf{Industry 4.0,}} often referred to as the \textit{fourth industrial revolution}, marks a significant paradigm shift in manufacturing and production processes. This revolution is characterized by the integration of digital technologies, such as the \textit{Internet of Things} (IoT), \textit{ Cyber-Physical systems } (CPS) \cite{baheti2011cyber,rose2015internet}, \textit{Big Data} and \textit{Artificial Intelligence} (AI), which together create highly interconnected and intelligent production environments. The aim of Industry 4.0 is to enhance the flexibility, efficiency, and sustainability of manufacturing processes, enabling businesses to adapt rapidly to changing market demands and technological advancements.




In \cite{hermann2016design}, the following primary design principles of Industry 4.0 are discussed.
\begin{enumerate}

 \item \textbf{Interconnection}: The ability of machines, devices, sensors, and people to connect and communicate with each other via the Internet of things.
 \item \textbf{Information transparency}: Access to comprehensive information to inform decisions. \item \textbf{Inter-connectivity}: Collect immense amounts of data and information from all points in the manufacturing process, identify key areas that can benefit from the improvement to increase functionality.
 \item \textbf{Technical assistance}: Modern technology to assist humans in decision-making and problem-solving, and the ability to help humans with difficult or unsafe tasks.
 \item \textbf{Decentralized decisions}: The ability of cyber-physical systems to make decisions on their own and to perform their tasks as autonomously as possible. Only in the case of exceptions, interference, or conflicting goals, are tasks delegated to a higher level
\end{enumerate}






